% !TeX root = ../../金融经济学_王江_讲义及习题.tex
% 知识点：references/金融经济学/金融经济学-王江_261013.pdf，印刷页 111–120（小结续至 121）
% 习题：references/金融经济学/史上最全最完整的金融经济学答案(王江版).pdf，页 32–36（印刷页 31–35）

\zhchapter{组合选择（一）}

\ifdefined\WJHandoutKnowledge

第 6、7 章给出了期望效用与风险厌恶的描述。本章在给定证券价格和支付的条件下，讨论参与者怎样选择消费与投资组合：先说明最优解为什么存在，再推导最优条件，最后分析风险溢价、风险厌恶与财富对投资的影响。具体效用函数下的显式解留到第 9 章。

\begin{tcolorbox}[knowledge,title=8.1 解的存在性及其特征]
\textbf{A.\ 消费与投资的联合选择}\quad 仍令 $X$ 为 $\Omega\times N$ 支付矩阵，$S$ 为价格列向量，$\theta$ 为证券持有\textbf{数量}。参与者的初始禀赋为 $(e_0,e_1)$，状态概率为 $\pi_\omega>0$，优化问题为
\[
\begin{aligned}
 \max_\theta\quad &u_0(c_0)+\sum_{\omega=1}^{\Omega}\pi_\omega u_1(c_{1\omega})\\
 \text{s.t.}\quad &c_0=e_0-S^{\mathsf T}\theta,\qquad c_1=e_1+X\theta,\\
 &c_0\ge0,\qquad c_1\ge0.
\end{aligned}\tag{8.1}
\]
投资占用今天的消费，却增加未来各状态下的消费。这里价格是参与者接受的外生条件，尚未讨论市场出清。

\begin{tcolorbox}[knowledge]
\textbf{定理 8.1}\quad 在有限状态、无摩擦且允许自由交易的模型中，若可行消费集合非空、效用连续并对每个时期及状态的消费严格递增，则消费与组合选择问题有最优解，当且仅当市场无套利。
\end{tcolorbox}

\textbf{证明（必要性）：}若 $\theta^*$ 最优，而 $\eta$ 是套利组合，则
\[
 d=\begin{bmatrix}-S^{\mathsf T}\eta\\X\eta\end{bmatrix}\ge0,\qquad d\ne0.
\]
选择 $\theta^*+\eta$ 后，消费由 $c^*$ 变为 $c^*+d$。没有任何分量减少，至少一项增加，故效用严格提高，与最优性矛盾。

\textbf{证明（充分性）：}无套利保证存在 $\phi\gg0$ 使 $S=X^{\mathsf T}\phi$。令 $\widehat\phi=(1,\phi^{\mathsf T})^{\mathsf T}$，任意可行消费满足
\[
 c_0+\phi^{\mathsf T}c_1=e_0+\phi^{\mathsf T}e_1=:W.
\]
由于消费非负且每个状态价格严格为正，
\[
 0\le c_0\le W,\qquad 0\le c_{1\omega}\le W/\phi_\omega.
\]
所以可行消费集合有界。它又是非负正交象限与仿射子空间
$e+\operatorname{Im}[-S^{\mathsf T};X]$ 的交集，因而闭。连续效用在这个非空有界闭集合上取得最大值，且该消费可由某个组合实现。

\textbf{说明：}有界的是可行\textbf{消费}，证券头寸未必有界；冗余证券允许不同头寸实现同一消费。严格凹效用保证最优消费唯一，组合数量还需支付映射无冗余才唯一。对数效用在零处无定义时，须在其定义域内处理边界；可利用趋近零时效用趋于 $-\infty$ 排除最优消费落在该边界。

\tcbline
\textbf{B.\ Euler 方程}\quad 把预算约束代入目标函数。若最优消费位于内点且效用可微，对第 $n$ 只证券的数量求导得
\[
 -u_0'(c_0)S_n+\sum_\omega\pi_\omega u_1'(c_{1\omega})x_{\omega,n}=0.
\]
所以
\[
 \underbrace{u_0'(c_0)S_n}_{\text{今天少消费的边际效用损失}}
 =\underbrace{\mathrm E[u_1'(\tilde c_1)\widetilde X_n]}_{\text{未来支付带来的期望边际效用}},
 \quad n=1,\ldots,N.\tag{8.2}
\]
这就是\textbf{Euler 方程}。当 $S_n\ne0$ 时，也可写为
\[
 \mathrm E\left[\frac{u_1'(\tilde c_1)}{u_0'(c_0)}
                 \frac{\widetilde X_n}{S_n}\right]=1.\tag{8.3}
\]
含义是：在最优处，将最后一元用于今天消费，或用于投资后在未来消费，两者的边际效用相等。令 $m_\omega=u_1'(c_{1\omega})/u_0'(c_0)$，则 $S_n=\mathrm E[m\widetilde X_n]$；$\pi_\omega m_\omega$ 给出由该参与者最优消费支持的一组状态价格。

\textbf{凹性为什么使一阶条件充分？}对任意头寸变动 $h$，目标函数的二阶方向导数为
\[
 u_0''(c_0)(S^{\mathsf T}h)^2+
 \sum_\omega\pi_\omega u_1''(c_{1\omega})(Xh)_\omega^2\le0.
\]
因此整个目标函数对 $\theta$ 凹，内点一阶条件给出全局最优。教材逐只证券的二阶条件为
\[
 S_n^2u_0''(c_0)+\mathrm E[u_1''(\tilde c_1)\widetilde X_n^2]\le0.\tag{8.4}
\]
多只证券时应同时看所有方向的凹性；只有各个对角项非正一般不足以判断最优。若消费非负、禁止卖空或借贷约束绑定，则应使用相应约束下的最优条件，不能直接套用内点等式。
\end{tcolorbox}

\begin{tcolorbox}[example]
\textbf{例 8.1（两状态下的对数效用）}\quad 两状态等概率，唯一证券价格为 $S>0$，支付为 $[x;1]$，其中 $x\ge1$。初始禀赋为 $e_0=w_0>0$、$e_1=0$，效用为
\[
 \log c_0+\frac{\rho}{2}(\log c_{1a}+\log c_{1b}),\qquad\rho>0.
\]
持有 $\theta$ 单位证券时，$c_0=w_0-S\theta$、$c_{1a}=\theta x$、$c_{1b}=\theta$，所以
\[
 \max_{0<\theta<w_0/S}\left\{\log(w_0-S\theta)+\rho\log\theta+
 \frac{\rho}{2}\log x\right\}.
\]
对 $\theta$ 求导，Euler 方程为
\[
 \frac{S}{w_0-S\theta}=\frac{\rho}{\theta}
 \ \Longrightarrow\ S\theta=\rho(w_0-S\theta)
 \ \Longrightarrow\ \theta^*=\frac{\rho}{1+\rho}\frac{w_0}{S}.
\]
二阶导数 $-S^2/(w_0-S\theta)^2-\rho/\theta^2<0$，且两端效用趋于 $-\infty$，故这是唯一最优解。相应消费为
\[
 c_0^*=\frac{w_0}{1+\rho},\qquad
 c_{1a}^*=\frac{\rho w_0x}{(1+\rho)S},\qquad
 c_{1b}^*=\frac{\rho w_0}{(1+\rho)S}.
\]
\textbf{经济含义：}投资支出 $S\theta^*=\rho w_0/(1+\rho)$；越重视未来（$\rho$ 越大），储蓄比例越高。财富增加使消费和投资同比增加；给定价格，提高 $x$ 只改变目标函数中的常数项，故不改变本例的最优数量。这是本例对数偏好与零未来禀赋的性质。

\textbf{核对边际效用：}
\[
 u_0'(c_0^*)=\frac{1+\rho}{w_0}
 =\frac{\rho}{2}\left(\frac{x}{Sc_{1a}^*}+\frac{1}{Sc_{1b}^*}\right).
\]
\end{tcolorbox}

\begin{tcolorbox}[knowledge,title=8.2 投资组合的选择]
\textbf{A.\ 先选组合，再选储蓄}\quad 令
\[
 w=e_0-c_0=S^{\mathsf T}\theta.\tag{8.5}
\]
这里 $w$ 是扣除当前消费后投入证券的\textbf{金融财富}。给定 $w$，最优组合能达到的期望效用称为\textbf{间接效用函数}：
\[
 v_1(w)=\max_{\theta:S^{\mathsf T}\theta=w}
       \mathrm E[u_1(\tilde e_1+\widetilde X\theta)].\tag{8.6}
\]
其中 $\widetilde X$ 是随机状态实现下的支付行向量，非负消费约束仍须满足。完整问题等价于
\[
 \max_w\{u_0(e_0-w)+v_1(w)\}.\tag{8.7}
\]
内层解决“同样的钱怎样投”，外层解决“今天消费多少、留下多少投资”。后文设 $e_1=0$，略去时间下标，重点研究
\[
 v(w)=\max_{\theta:S^{\mathsf T}\theta=w}\mathrm E[u(\widetilde X\theta)].\tag{8.8}
\]
不可交易的未来收入会改变组合选择，不能直接并入可自由投资的 $w$。

\tcbline
\textbf{B.\ 从证券数量改写成投资金额}\quad 设第 $N$ 只证券无风险，净利率为 $r_F$，$1+r_F>0$。将证券取向和计量单位适当调整后，设 $S_n>0$，定义
\[
 a_n=S_n\theta_n,\qquad
 \tilde x_n=\frac{\widetilde X_n}{S_n},\tag{8.9}
\]
\[
 \tilde r_n=\frac{\widetilde X_n-S_n}{S_n}=\tilde x_n-1.\tag{8.10}
\]
$a_n$ 是\textbf{金额}，$\theta_n$ 是\textbf{数量}，$a_n/w$ 才是投资\textbf{权重}；$\tilde x_n$ 是总收益率，$\tilde r_n$ 是净收益率。投资金额满足 $\sum_n a_n=w$，从而期末财富为
\[
\begin{aligned}
 \tilde w&=\sum_{n=1}^N a_n(1+\tilde r_n)\\
 &=w(1+r_F)+\sum_{n=1}^{N-1}a_n(\tilde r_n-r_F).
\end{aligned}\tag{8.11}
\]
$\tilde r_n-r_F$ 是\textbf{超额收益率}，其期望 $\mathrm E[\tilde r_n]-r_F$ 是本章的\textbf{风险溢价}（收益率之差），与第 7 章用确定性等价定义的货币风险溢价须区分。

记 $a=(a_1,\ldots,a_{N-1})^{\mathsf T}$，$\iota$ 为全 1 列向量，组合问题变成
\[
 \max_a\ \mathrm E\bigl[u(w(1+r_F)+a^{\mathsf T}(\tilde r-r_F\iota))\bigr].\tag{8.12}
\]
无风险投资为 $a_N=w-\sum_{n<N}a_n$；负数表示借款。无约束内点解满足
\[
 \mathrm E[u'(\tilde w)(\tilde r_n-r_F)]=0,
 \qquad n=1,\ldots,N-1.\tag{8.13}
\]
这表示在风险证券与无风险证券之间再转移一元，期望边际效用的净变化为零。
\end{tcolorbox}

\begin{tcolorbox}[knowledge,title=8.3 最优组合的性质：单只风险证券]
以下设 $u'>0$、$u''<0$，风险证券收益非退化，允许卖空与借贷，并假定有限的内点最优解存在、期望与求导可交换。令
\[
 z=\tilde r-r_F,\qquad b=w(1+r_F),\qquad
 F(a)=\mathrm E[u(b+az)].
\]
则 $F'(a)=\mathrm E[u'(b+az)z]$，$F''(a)=\mathrm E[u''(b+az)z^2]<0$。

\begin{tcolorbox}[knowledge]
\textbf{定理 8.2}\quad 最优风险投资额 $a^*$ 与风险溢价同号：
\[
 \begin{array}{c|ccc}
 \mathrm E[\tilde r]-r_F&=0&>0&<0\\\hline
 a^*&=0&>0&<0
 \end{array}
\]
\end{tcolorbox}
\textbf{证明：}在 $a=0$ 时，$F'(0)=u'(b)\mathrm E[z]$，其符号与 $\mathrm E[z]$ 相同。严格凹性使 $F'$ 严格递减，故 $F'(a^*)=0$ 的根分别在零点、零点右侧或左侧。正风险溢价使参与者愿意从零开始买入，但持仓增大时风险成本最终抵消收益。

\tcbline
\textbf{A.\ 绝对风险厌恶决定投资金额怎样随财富变化}\quad 以下设 $\mathrm E[z]>0$，所以 $a^*(w)>0$。记
$A(t)=-u''(t)/u'(t)$。
\begin{tcolorbox}[knowledge]
\textbf{定理 8.3}\quad 在相关财富区间内：CARA 使最优风险投资\textbf{金额}不随财富变化；DARA 使其随财富增加；IARA 使其随财富减少。
\end{tcolorbox}
\textbf{证明：}把一阶条件 $\mathrm E[u'(b+a^*z)z]=0$ 对 $w$ 求导，记 $\tilde w^*=b+a^*z$，得
\[
 (1+r_F)\mathrm E[u''(\tilde w^*)z]
 +\frac{da^*}{dw}\mathrm E[u''(\tilde w^*)z^2]=0,
\]
即
\[
 \frac{da^*}{dw}=-(1+r_F)
 \frac{\mathrm E[u''(\tilde w^*)z]}{\mathrm E[u''(\tilde w^*)z^2]}.
\]
分母负，故导数与分子 $\mathrm E[u''(\tilde w^*)z]$ 同号。利用一阶条件减掉常数项，
\[
 \mathrm E[u''(\tilde w^*)z]
 =-\mathrm E[(A(\tilde w^*)-A(b))u'(\tilde w^*)z].
\]
DARA 下，$z>0$ 时 $\tilde w^*>b$、$A(\tilde w^*)<A(b)$；$z<0$ 时则反向。两种情况下括号与 $z$ 的乘积均为负，前面再取负号，得到 $da^*/dw>0$。CARA 时差为零；IARA 时符号相反。

这里比较的是 $A$ 在可能期末财富\textbf{整个区间}上的变化，不能仅凭初始财富处一个 $A'(w)$ 的符号推断任意风险下的结论。

\tcbline
\textbf{B.\ 相对风险厌恶决定投资比例怎样随财富变化}\quad 设 $w>0$、期末财富为正，令
\[
 \alpha^*(w)=\frac{a^*(w)}{w},\qquad
 e(w)=\frac{w}{a^*(w)}\frac{da^*}{dw}.
\]
$e$ 是风险投资金额对财富的\textbf{弹性}，且
\[
 \frac{d\alpha^*}{dw}=\frac{a^*}{w^2}(e-1).
\]
\begin{tcolorbox}[knowledge]
\textbf{定理 8.4}\quad 记 $R(t)=tA(t)$。在相关财富区间内，CRRA 对应 $e=1$、风险投资比例不变；DRRA 对应 $e>1$、比例随财富增加；IRRA 对应 $e<1$、比例随财富减少。
\end{tcolorbox}
\textbf{证明：}记 $D=\mathrm E[u''(\tilde w^*)z^2]<0$。由前式及 $\tilde w^*=b+a^*z$，
\[
\begin{aligned}
 e-1&=-\frac{b\,\mathrm E[u''(\tilde w^*)z]+a^*D}{a^*D}\\
 &=-\frac{\mathrm E[u''(\tilde w^*)\tilde w^*z]}{a^*D}
 =\frac{\mathrm E[R(\tilde w^*)u'(\tilde w^*)z]}{a^*D}.
\end{aligned}
\]
再用一阶条件，将分子写成
$\mathrm E[(R(\tilde w^*)-R(b))u'(\tilde w^*)z]$。
DRRA 下该分子为负，分母也为负，故 $e-1>0$；CRRA 时为零，IRRA 时相反。\textbf{金额增加不等于比例增加：}例如 CRRA 下金额随财富同比增加，而比例不变。
\end{tcolorbox}

\begin{tcolorbox}[knowledge,title=8.3（续）最优组合的性质：多只风险证券]
\begin{tcolorbox}[knowledge]
\textbf{定理 8.5}\quad 所有风险证券的风险溢价都为零，当且仅当全投无风险证券的 $a=0$ 是最优组合之一。
\end{tcolorbox}
\textbf{证明：}若 $a=0$ 最优，内点一阶条件给出
$u'(b)\mathrm E[\tilde r_n-r_F]=0$，故每只证券风险溢价均为零。反之，若均为零，则任意组合都有 $\mathrm E[\tilde w]=b$。由 Jensen 不等式，
\[
 \mathrm E[u(\tilde w)]\le u(\mathrm E[\tilde w])=u(b),
\]
而 $a=0$ 恰好达到 $u(b)$，故最优。若存在能完全抵消的风险头寸，最优金额向量未必唯一；严格凹性排除的是带有非零消费风险的同效用方案。

\begin{tcolorbox}[knowledge]
\textbf{定理 8.6}\quad 最优组合的期望收益率不低于无风险收益率。若至少一只证券风险溢价非零，且允许在零头寸附近自由调整、最优解存在，则最优组合的期望收益率严格更高。
\end{tcolorbox}
\textbf{证明：}全投无风险证券是可行的，所以
\[
 u(\mathrm E[\tilde w^*])\ge\mathrm E[u(\tilde w^*)]\ge u(b).
\]
由 $u'>0$，$\mathrm E[\tilde w^*]\ge b$，即
\[
 (a^*)^{\mathsf T}(\mathrm E[\tilde r]-r_F\iota)\ge0.
\]
若某一风险溢价非零，在 $a=0$ 处沿该证券风险溢价的同号方向小幅交易，效用的一阶变化严格为正，所以最优效用严格大于 $u(b)$，上述期望财富不等式也严格成立。$w>0$ 时除以 $w$，得到期望组合净收益率大于 $r_F$。

\textbf{为什么不能逐只套用定理 8.2？}多资产下还要看证券对\textbf{整个组合风险}的影响。即使每只证券的风险溢价都为正，最优组合也可能卖空其中某只证券，以对冲其余证券的风险；只能保证最优组合整体的期望超额收益非负。
\end{tcolorbox}

\begin{tcolorbox}[knowledge,title=8.4 本章结论]
\begin{enumerate}
 \item 无套利使可行消费集合有界；连续效用在非空闭集合上达到最优。冗余会影响头寸唯一性。
 \item 内点 Euler 方程把今天投资的边际成本与未来消费的期望边际效用相等；凹性使一阶条件充分。
 \item 消费与投资可分两步：给定储蓄求最优组合，再用间接效用决定储蓄。
 \item 单风险资产下，最优头寸与风险溢价同号；绝对风险厌恶控制投资\textbf{金额}的财富效应，相对风险厌恶控制投资\textbf{比例}的财富效应。
 \item 多风险资产下，零风险溢价使全投无风险证券最优；存在非零风险溢价时，最优组合整体有正的期望超额收益，但不能据此推断每只证券的持仓符号。
\end{enumerate}
\end{tcolorbox}

\fi

% !TeX root = ../../金融经济学_王江_讲义及习题.tex
% 08_组合选择一 课后习题（由章节文件拆分）

\begin{problem}
\problemno{8.1} 假设有 $N$ 项资产，收益率为 $\tilde{r}_{i}$，$i=1,\cdots,N$，它们是独立同分布的。对于一个具有不满足性和凹性效用函数的投资者，计算他的最优投资组合。
\end{problem}
\begin{solution}
设组合权重 $\theta=(\theta_{1},\ldots,\theta_{N})$，$\sum_{i}\theta_{i}=1$，期望效用目标先写为
\[
 \max_\theta E\!\left[u\!\left(
 \underbrace{w\sum_i\theta_i\tilde r_i}_{\tilde w\text{：期末财富}}
 \right)\right].
\]
其一阶条件为：对任意 $j,k$，
\[
    E\!\left[u'\!\left(w\sum_i\theta_i\tilde r_i\right)(\tilde r_j-\tilde r_k)\right]=0.
\]
由于 $\{\tilde{r}_{i}\}$ 独立同分布，目标函数对 $\theta$ 对称；严格凹效用下最优解唯一，故必有 $\theta_{i}^{*}=1/N$。即最优组合是所有风险资产的等权重组合。
\end{solution}

\begin{problem}
\problemno{8.2*} 假设你必须将 $w$ 投资于两只证券：一只无风险债券和一只风险证券。无风险债券的收益率确定为 $r_{F}$。风险证券的收益率为 $\tilde{r}$，它的均值为 $\bar{r}$ 而方差为 $\sigma^{2}$。假设你的目标就是最大化你的二次效用函数的期望值：
\[
    E\Bigl[\tilde{w}-\tfrac12 a\tilde{w}^{2}\Bigr].
\]
\begin{enumerate}[label=(\alph*)]
    \item 求解最优组合。
    \item 讨论最优组合如何依赖于股票和债券的期望收益率 $r_{F}$ 和 $\bar{r}$，股票收益率的波动率 $\sigma$，以及偏好系数 $a$。为这些依赖关系提供经济解释。
\end{enumerate}
\end{problem}
\begin{solution}
\begin{enumerate}[label=(\alph*)]
    \item 假设投资于股票的资产比例为 $x$，那么有
    \[
        \tilde{w}=w\bigl(1+r_{F}+x(\tilde{r}-r_{F})\bigr)
    \]
    参与者的优化问题是
    \[
        \max_x E\!\left[
 \underbrace{w\bigl(1+r_F+x(\tilde r-r_F)\bigr)}_{\tilde w}
 -\frac a2\bigl(w(1+r_F+x(\tilde r-r_F))\bigr)^2\right].
    \]
    先计算期末财富的一阶、二阶矩：
    \[
    \begin{aligned}
        E[\tilde w]&=w\bigl(1+r_F+x(\bar r-r_F)\bigr),\\
        E[\tilde w^2]&=w^2\bigl[(1+r_F+x(\bar r-r_F))^2+x^2\sigma^2\bigr].
    \end{aligned}
    \]
    代入目标函数后对 $x$ 求导，一阶条件为
    \[
        w(\bar r-r_F)-aw^2\bigl[(1+r_F+x(\bar r-r_F))(\bar r-r_F)+x\sigma^2\bigr]=0.
    \]
    整理得
    \[
        (\bar r-r_F)\bigl[1-aw(1+r_F)\bigr]
        =aw\bigl((\bar r-r_F)^2+\sigma^2\bigr)x.
    \]
    解得最优权重
    \[
        x^{*}=\frac{(\bar{r}-r_{F})\bigl(1-aw(1+r_{F})\bigr)}{aw\bigl(\sigma^{2}+(\bar{r}-r_{F})^{2}\bigr)}
    \]

    \item $x^{*}$ 随着 $\bar{r}$ 的增加而增加，随着 $r_{F}$、$\sigma^{2}$、$a$ 的增加而减小。股票的期望收益率增加使得股票的需求增加，而无风险收益率增加使得债券的需求增加；$\sigma^{2}$ 增大，股票的风险变大，对股票的需求变小；最后，$a$ 增大表明参与者更加厌恶风险，因而对股票的需求减小。
\end{enumerate}
\end{solution}

\begin{problem}
\problemno{8.3*} 考虑与练习 8.2 一样的问题。不过现在假设效用函数具有常数绝对风险厌恶系数：
\[
    E\bigl[-e^{-a\tilde{w}}\bigr].
\]
另外，假设股票收益率 $\tilde{r}$ 是正态分布的。
\begin{enumerate}[label=(\alph*)]
    \item 求解最优投资组合。
    \item 讨论最优组合如何依赖于股票和债券的期望收益率 $r_{F}$ 和 $\bar{r}$，股票收益率的波动率 $\sigma$，以及偏好系数 $a$。为这些依赖性提供经济解释。[如果需要的话，你可以添加其它条件。]
\end{enumerate}
\end{problem}
\begin{solution}
\begin{enumerate}[label=(\alph*)]
    \item 假设投资于股票的资产比例为 $x$，那么有
    \[
        \tilde{w}=w\bigl(1+r_{F}+x(\tilde{r}-r_{F})\bigr)
    \]
    参与者的优化问题是
    \[
        \max_x E\!\left[-\exp\!\left(-a
 \underbrace{w(1+r_F+x(\tilde r-r_F))}_{\tilde w}\right)\right].
    \]
    因 $\tilde{r}$ 正态，$\tilde{w}$ 亦正态。对正态随机变量有
    \[
        E\bigl[-e^{-a\tilde{w}}\bigr]=-\exp\Bigl(-aE[\tilde{w}]+\tfrac12 a^{2}\operatorname{Var}(\tilde{w})\Bigr)
    \]
    其中 $E[\tilde{w}]=w\bigl(1+r_{F}+x(\bar{r}-r_{F})\bigr)$，$\operatorname{Var}(\tilde{w})=w^{2}x^{2}\sigma^{2}$。最大化上式等价于最大化确定性等价，完整写为
    \[
        \max_x\left\{w\bigl(1+r_F+x(\bar r-r_F)\bigr)
        -\frac a2 w^2x^2\sigma^2\right\}.
    \]
    对 $x$ 的一阶条件为
    \[
        w(\bar{r}-r_{F})-a\,w^{2}x\sigma^{2}=0
    \]
    因而最优权重（CARA 下风险资产的金额 $wx^{*}$ 与财富无关）为
    \[
        x^{*}=\frac{\bar{r}-r_{F}}{aw\sigma^{2}}
    \]
    \item $x^{*}$ 随着 $\bar{r}$ 的增加而增加，随着 $r_{F}$、$\sigma^{2}$、$a$ 以及财富 $w$ 的增加而减小。超额收益上升提高股票需求；$\sigma^{2}$ 或绝对风险厌恶 $a$ 增大则降低股票需求。与二次效用不同：CARA 下风险资产金额 $wx^{*}=(\bar{r}-r_{F})/(a\sigma^{2})$ 不随 $w$ 变化，故权重 $x^{*}$ 随财富下降。
\end{enumerate}
\end{solution}

\begin{problem}
\problemno{8.4**} 赋予投资者初始金融财富 $w$ 以及 $\lambda$ 份不可交易收入 $\tilde{y}$，其中 $\lambda\ge 0$ 且 $E[\tilde{y}]=0$。他可以交易一只无风险证券和一只风险证券。无风险证券的收益率为 0 而风险证券的收益率 $\tilde{r}$ 的均值为 $\bar{r}$。他具有不满足性且是风险厌恶的。他的绝对风险厌恶系数随着他的财富增加而减少。令 $a$ 为风险证券上的投资额。他的期望效用函数为：
\[
    \max E\bigl[u(w+a\tilde{r}+\lambda\tilde{y})\bigr].
\]
\begin{enumerate}[label=(\alph*)]
    \item 当 $\lambda=0$ 时，风险资产的最优头寸 $a_{0}^{*}$ 如何随着 $w$ 变化而变化？
    \item 假设 $\tilde{y}$ 和 $\tilde{r}$ 是独立的。当 $\lambda>0$ 时，记风险资产的最优头寸为 $a_{\lambda}^{*}$。$a_{\lambda}^{*}$ 是大于还是小于 $a_{0}^{*}$？
    \item 假设 $\tilde{y}=c(\tilde{r}-\bar{r})+\tilde{\varepsilon}$ 以及 $E[\tilde{\varepsilon}\mid\tilde{r}]=0$。因此，如果 $c\neq 0$，$\tilde{y}$ 和 $\tilde{r}$ 是相关的。当 $\lambda$ 极小时，他持有的股票头寸如何随着 $w$ 变化而变化？解释所得到的结论。
\end{enumerate}
\end{problem}
\begin{solution}
\begin{enumerate}[label=(\alph*)]
    \item 当 $\lambda=0$ 时，优化的一阶条件是
    \[
        E\bigl[u'\!\left(\underbrace{w+a_0^*\tilde r}_{\tilde w}\right)\tilde{r}\bigr]=0
    \]
    两边对 $w$ 求导可得
    \[
        \frac{da_{0}^{*}}{dw}=\frac{E[u''(\tilde{w})\tilde{r}]}{-E[u''(\tilde{w})\tilde{r}^{2}]}
    \]
    这里绝对风险厌恶系数沿用正文的记号，写为
    \[
        \underbrace{-\frac{u''(t)}{u'(t)}}_{A(t)\text{：绝对风险厌恶系数}}.
    \]
    又
    \[
        \begin{aligned}
            E[u''(\tilde{w})\tilde{r}]
            &=-E[u'(\tilde{w})A(\tilde{w})\tilde{r}]\\
            &\ge -E[u'(\tilde{w})A(w)\tilde{r}]\\
            &=-A(w)E[u'(\tilde{w})\tilde{r}]\\
            &=0
        \end{aligned}
    \]
    故 $\dfrac{da_{0}^{*}}{dw}\ge 0$，即最优头寸随着 $w$ 的增加而增加；

    \item 参与者优化的一阶条件是
    \[
        E\bigl[u'\!\left(\underbrace{w+a_\lambda^*\tilde r+\lambda\tilde y}_{\tilde w}\right)\tilde{r}\bigr]=0
    \]
    两边对 $\lambda$ 求导可得
    \[
        \frac{da_{\lambda}^{*}}{d\lambda}=\frac{E[u''(\tilde{w})\tilde{r}\tilde{y}]}{-E[u''(\tilde{w})\tilde{r}^{2}]}
    \]
    将下面的条件期望记为 $K$：
    \[
        \underbrace{E[-A(\tilde w)\tilde y\mid\tilde r]}_{K(w+a_\lambda^*\tilde r)}.
    \]
    可以证明 $K(\cdot)>0$，$K'(\cdot)>0$，因此
    \[
        \begin{aligned}
 E[u''(\tilde w)\tilde r\tilde y]
 &=E[-A(\tilde w)u'(\tilde w)\tilde r\tilde y]\\
 &=E[K(w+a_\lambda^*\tilde r)u'(\tilde w)\tilde r]\\
 &\ge K(w)E[u'(\tilde w)\tilde r]=0.
 \end{aligned}
    \]
    最后一个等式是因为一阶条件，因此，$\dfrac{da_{\lambda}^{*}}{d\lambda}\ge 0$，故 $a_{\lambda}^{*}\ge a_{0}^{*}$；

    \item 由 (a) 有
    \[
        \frac{da_{0}^{*}}{dw}=\frac{E[u''(\tilde{w})\tilde{r}]}{-E[u''(\tilde{w})\tilde{r}^{2}]}=\frac{-E[u'(\tilde{w})A(\tilde{w})\tilde{r}]}{-E[u''(\tilde{w})\tilde{r}^{2}]}
    \]
    而
    \[
        \begin{aligned}
 E[u'(\tilde w)A(\tilde w)\tilde r]
 &=E[A(w+a\tilde r)u'(\tilde w)\tilde r]\\
 &\quad+\lambda E[A'(w+a\tilde r)u'(\tilde w)\tilde r\tilde y]
 +o(\lambda^2).
 \end{aligned}
    \]
    由 (a) 知道 $E\bigl[A(w+a\tilde{r})u'(\tilde{w})\tilde{r}\bigr]\le 0$，另外 $\lambda$ 又极小，因此
    \[
        E[u'(\tilde{w})A(\tilde{w})\tilde{r}]\le 0,\qquad \frac{da^{*}}{dw}\ge 0
    \]
    即最优头寸随着初始财富 $w$ 的增加而增加。这是因为 $\lambda$ 极小，因而不可交易收入不会改变 (a) 中的结论。
\end{enumerate}
\end{solution}

\begin{problem}
\problemno{8.5**} 假设你是对冲基金中交易政府债的交易员。你有 $w$ 的现金额度。
\begin{itemize}
    \item 在 0 期，当美联储拍卖新债券时，你可以提交报价而买入政府债。在 1 期，你可以在二级市场上以价格 $\tilde{v}$ 再把它们卖掉，其中 $\tilde{v}=\bar{v}+\tilde{\varepsilon}$，$\bar{v}>0$ 且 $E[\tilde{\varepsilon}]=0$。假设无风险收益率为零。
    \item 你购买债券的报价是 $p$。你购买到的数量 $x(p)$ 是 $p$ 的递增函数且 $x(0)=0$。选择 $p$ 以最大化你的期望效用函数：
    \[
        V(p)=E[u(\tilde{w})],\qquad \tilde{w}=w+x(p)(\bar{v}+\tilde{\varepsilon}-p).
    \]
    假设你是不满足的且严格风险厌恶的。为简单起见，假设 $V(p)$ 是严格凹的。
\end{itemize}
\begin{enumerate}[label=(\alph*)]
    \item 写出关于最优报价 $p^{*}$ 的一阶条件。
    \item 证明 $0<p^{*}<\bar{v}$。
    \item 对于两个绝对风险厌恶系数分别为 $A_{1}(w)$ 和 $A_{2}(w)$ 的交易员，如果 $\forall\,w$，$A_{1}(w)\ge A_{2}(w)$，那么 $p_{1}^{*}\le p_{2}^{*}$，其中 $p_{i}^{*}$ 是交易员 $i$ 的最优报价，$i=1,2$。
    \item 假设 $u(\cdot)$ 是绝对风险厌恶递减的。证明初始财富 $w$ 的增加使得 $p^{*}$ 也增加。
\end{enumerate}
\end{problem}
\begin{solution}
\begin{enumerate}[label=(\alph*)]
    \item 一阶条件是
    \[
        \begin{aligned}
 V'(p^*)&=E\!\left[
 u'\!\left(\underbrace{w+x(p^*)(\bar v+\tilde\varepsilon-p^*)}_{\tilde w}\right)\right.\\
 &\hspace{3em}\left.\cdot\bigl(x'(p^*)(\bar v+\tilde\varepsilon-p^*)-x(p^*)\bigr)\right]=0.
 \end{aligned}
    \]

    \item
    \[
        \begin{aligned}
            V'(0)&=E\bigl[u'(w)\bigl(x'(0)(\bar{v}+\tilde{\varepsilon})\bigr)\bigr]=u'(w)x'(0)\bar{v}>0\\
            V(\bar{v})&=E\bigl[u(w+x(\bar{v})\tilde{\varepsilon})\bigr]<u(w)=V(0)
        \end{aligned}
    \]
    又因为 $V(p)$ 是凹的，故 $0<p^{*}<\bar{v}$；

    \item 由定理 (7.4) 可知，存在 $f'>0$，$f''<0$ 使得 $u_{1}=f(u_{2})$。参与者的优化问题的一阶条件为：
    \[
        \begin{aligned}
 E\!\left[
 u_2'\!\left(\underbrace{w+x(p_2)(\bar v+\tilde\varepsilon-p_2)}_{\tilde w_2}\right)\right.&\\
 \left.\cdot\bigl(x'(p_2)(\bar v+\tilde\varepsilon-p_2)-x(p_2)\bigr)\right]&=0.
 \end{aligned}
    \]
    考虑
    \[
        V_{1}'(p_{2})=E\bigl[f'(u_{2}(\tilde{w}_{2}))u_{2}'(\tilde{w}_{2})\bigl(x'(p_{2})(\bar{v}+\tilde{\varepsilon}-p_{2})-x(p_{2})\bigr)\bigr]
    \]
    将报价的边际财富变化标记为
    \[
        \underbrace{x'(p_2)(\bar v+\tilde\varepsilon-p_2)-x(p_2)}_{\tilde r}.
    \]
    那么 $\tilde{w}_{2}$ 是 $\tilde{r}$ 的增函数，又 $f'(u_{2}(\tilde{w}_{2}))>0$ 且是 $\tilde{w}_{2}$ 的减函数，因此，$V_{1}'(p_{2})\le 0$，另外 $V(p)$ 是凹函数，因此 $p_{1}^{*}\le p_{2}^{*}$。

    \item 假设 $w_{1}>w_{0}$，当 $w=w_{0}$ 时优化问题的一阶条件为
    \[
        \begin{aligned}
 E\!\left[u'\!\left(\underbrace{w_0+x(p_0)(\bar v+\tilde\varepsilon-p_0)}_{\tilde w_0}\right)\right.&\\
 \left.\cdot\bigl(x'(p_0)(\bar v+\tilde\varepsilon-p_0)-x(p_0)\bigr)\right]&=0.
 \end{aligned}
    \]
    以下仍将绝对风险厌恶系数写为
    \[
        \underbrace{-\frac{u''(t)}{u'(t)}}_{A(t)}.
    \]
    等式两边对 $w_{0}$ 求导可得
    \[
        \begin{aligned}
            &E\bigl[u''(\tilde{w}_{0})\bigl(x'(p_{0})(\bar{v}+\tilde{\varepsilon}-p_{0})-x(p_{0})\bigr)\bigr]\\
            &=E\bigl[-A(\tilde{w}_{0})u'(\tilde{w}_{0})\bigl(x'(p_{0})(\bar{v}+\tilde{\varepsilon}-p_{0})-x(p_{0})\bigr)\bigr]\\
            &\ge 0
        \end{aligned}
    \]
    因此，
    \[
        E\bigl[u'\bigl(w_{1}+x(p_{0})(\bar{v}+\tilde{\varepsilon}-p_{0})\bigr)\bigl(x'(p_{0})(\bar{v}+\tilde{\varepsilon}-p_{0})-x(p_{0})\bigr)\bigr]\ge 0
    \]
    而 $V(p)$ 又是凹的，因此 $p_{1}\ge p_{0}$，即初始财富增加报价也会增加。
\end{enumerate}
\end{solution}


